% ============================================================================
% MODULE 4: NITROGEN FERTILISATION
% ============================================================================
% Duration: 60 minutes (N response, AEN, optimization)
% Learning Goals: N response curves; AEN calculation; cost-benefit analysis
% Prerequisites: Module 3

\chapter{Module 4: Nitrogen Fertilisation}
\label{ch:module4}

\section{Learning Objectives}

By the end of this module, you will:
\begin{itemize}
  \item Understand N cycling in APSIM (soil N pools, uptake, plant demand)
  \item Generate N response curves (yield vs. N rate)
  \item Calculate AEN (Agronomic Efficiency of Nitrogen)
  \item Assess N stress effects on crop development and protein
  \item Optimize N rate based on cost-benefit analysis
\end{itemize}

\section{Nitrogen Cycling in APSIM}

\subsection{Soil N Pools}

\begin{itemize}
  \item \textbf{NO\textsubscript{3} (nitrate):} Mobile; easily leached; available to plant
  \item \textbf{NH\textsubscript{4} (ammonium):} Less mobile; adsorbed to clay; slowly nitrified
  \item \textbf{Organic N:} Immobilized in soil organic matter; mineralized slowly
  \item \textbf{Urea:} Applied as fertiliser; hydrolyzed to NH\textsubscript{4} over days
\end{itemize}

\subsection{Daily N Processes}

\begin{enumerate}
  \item \textbf{Mineralization:} Organic N $\to$ NH\textsubscript{4}
  \item \textbf{Nitrification:} NH\textsubscript{4} $\to$ NO\textsubscript{3}
  \item \textbf{Plant uptake:} N availability $\times$ plant demand $=$ actual uptake
  \item \textbf{Leaching:} NO\textsubscript{3} lost with drainage water
  \item \textbf{Immobilization:} N incorporated into new microbial biomass
\end{enumerate}

\section{N Response Curves}

\subsection{Typical Response}

When N rate increases from 0 to a high rate, yield follows a diminishing-returns curve:

\begin{enumerate}
  \item At \textbf{low N} (0--80 kg/ha): yield increases steeply — N is strongly limiting
  \item At \textbf{moderate N} (80--160 kg/ha): yield still increases but more slowly
  \item At \textbf{high N} ($>$160 kg/ha): yield approaches a plateau — other factors limit (water, light, genetics)
\end{enumerate}

Typical wheat N response at Northam (Red/Brown earth, Mediterranean WA Wheatbelt):

\begin{center}
\begin{tabular}{r|r|r|r}
\textbf{N Rate (kg/ha)} & \textbf{Yield (kg/ha)} & \textbf{Yield gain (kg/ha)} & \textbf{NUE (kg/kg)} \\
\hline
0   & 1800 & ---  & --- \\
80  & 2960 & 1160 & 14.5 \\
160 & 3780 & 820  & 12.4 \\
240 & 4060 & 280  & 9.4 \\
\end{tabular}
\end{center}

\begin{tipbox}
\textbf{Note:} The values above are indicative. Your simulated values will differ depending on the season's rainfall and initial soil N. In Exercise B you will calculate these numbers from your own simulation output.
\end{tipbox}

% ---- IMAGE PLACEHOLDER ----
\begin{figure}[H]
  \centering
  \fbox{\parbox{0.85\textwidth}{\centering\vspace{1.5cm}
    \textit{[IMAGE: N response curve plot. X-axis: N rate (0, 80, 160, 240 kg/ha). Y-axis: Grain yield (kg/ha).}\\
    \textit{Show diminishing returns shape. Mark optimal economic N rate with vertical dashed line.]}
  \vspace{1.5cm}}}
  \caption{Typical wheat N response curve showing diminishing returns at high N rates. The optimal N rate depends on grain price and fertiliser cost.}
  \label{fig:n-response-curve}
\end{figure}
% ---- END PLACEHOLDER ----

\subsection{Exercise B: N Response Factorial}

Four N rate variants:
\begin{itemize}
  \item 0 kg/ha (control)
  \item 80 kg/ha
  \item 160 kg/ha (standard)
  \item 240 kg/ha (high)
\end{itemize}

Run all four; compare yields, grain protein, and AEN.

\section{Agronomic Efficiency of Nitrogen}

\[
\text{AEN} = \frac{\text{Yield}_{N \text{ treatment}} - \text{Yield}_{0 \text{ N}}}{\text{N rate}} \quad [\text{kg grain / kg N}]
\]

\begin{keyconceptbox}
\textbf{Interpretation:}
\begin{itemize}
  \item AEN = 15 kg/kg: Every kg of N produces 15 kg of grain
  \item AEN is typically highest at low N rates (diminishing returns)
  \item AEN decreases as N rate increases (physiological limit)
\end{itemize}
\end{keyconceptbox}

\section{N Stress and Crop Responses}

When soil N supply cannot meet crop demand, APSIM applies an N stress factor that reduces growth:

\begin{center}
\begin{tabular}{l|l|l}
\textbf{Stressed component} & \textbf{Effect} & \textbf{Visible in output} \\
\hline
Leaf area (LAI) & Reduced leaf expansion & Lower LAI peak \\
Biomass accumulation & Slower carbon fixation & Lower biomass at anthesis \\
Grain number & Fewer florets set & Lower grain count at Stage 65 \\
Grain protein & Higher protein if grain N $>$ carbohydrate & \variable{[Wheat].GrainProtein} rises \\
Harvest Index & Grain fill incomplete & HI $<$ 0.40 \\
\end{tabular}
\end{center}

\begin{keyconceptbox}
\textbf{N stress vs water stress:} Both reduce yield, but through different mechanisms.
\begin{itemize}
  \item \textbf{N stress}: Reduces leaf area and biomass accumulation (fewer sources of carbon)
  \item \textbf{Water stress}: Reduces stomatal conductance and transpiration (fewer sinks fulfilled)
  \item In APSIM, the more limiting stress dominates on any given day
\end{itemize}
\end{keyconceptbox}

\section{Cost-Benefit Analysis}

\subsection{Profit Formula}

\[
\text{Net Profit (\$/ha)} = (\text{Yield (kg/ha)} \times \text{Grain price (\$/kg)}) - (\text{N rate (kg/ha)} \times \text{N cost (\$/kg N)})
\]

Typical values for Queensland dryland wheat (indicative only):
\begin{itemize}
  \item Grain price: \$0.28--0.35 per kg (\$280--350 per tonne)
  \item Urea (46\% N): approximately \$1.20--1.60 per kg N (varies with market)
  \item Application cost: approximately \$15--25 per ha
\end{itemize}

\subsection{Worked Example}

Using the indicative N response values from Section 4.3:

\begin{center}
\begin{tabular}{r|r|r|r|r}
\textbf{N Rate} & \textbf{Yield} & \textbf{Revenue} & \textbf{N Cost} & \textbf{Net Profit} \\
(kg/ha) & (kg/ha) & @ \$0.30/kg & @ \$1.40/kg N & (\$/ha) \\
\hline
0   & 1800 & 540  & 0   & 540 \\
80  & 2960 & 888  & 112 & 776 \\
160 & 3780 & 1134 & 224 & 910 \\
240 & 4060 & 1218 & 336 & 882 \\
\end{tabular}
\end{center}

In this example, 160 kg/ha gives the highest net profit. At 240 kg/ha, the additional yield does not justify the extra fertiliser cost.

\begin{warningbox}
\textbf{Year-to-year variability matters.} In a dry year, yield potential is lower, so the optimal N rate is lower. In Exercise B you will compare profit across a single season. Exercise D extends this to 20 years to reveal how profit varies with climate.
\end{warningbox}

\section{Regional Variability and AEN Interpretation}

AEN varies significantly across regions because it depends on soil N mineralisation, rainfall, and variety:

\begin{center}
\begin{tabular}{l|l|l}
\textbf{Region type} & \textbf{Typical AEN (kg grain/kg N)} & \textbf{Key driver} \\
\hline
High-rainfall, leaching soil & 8--12 & N lost to drainage; moderate response \\
Subtropical clay soil (Northam, WA) & 8--12 & Variable N retention; moderate response \\
Drought-prone, low rainfall & 6--10 & Water limits yield regardless of N \\
High-mineralising soil & 8--12 & Soil supplies some N; crop less responsive \\
\end{tabular}
\end{center}

\begin{keyconceptbox}
\textbf{Why WA Red/Brown earth soils respond to N:}
The Red/Brown earth soil at Northam has moderate cation-exchange capacity, which adsorbs NH$_4^+$ and reduces leaching. However, it drains more freely than Black Vertosols, so N management requires attention to timing and rainfall. NUE on these soils is typically moderate (10--15 kg/kg), reflecting both the soil's capacity to retain applied N and the constraints of the Mediterranean climate.
\end{keyconceptbox}

\section{Summary}

\begin{keyconceptbox}
\textbf{Module 4 key points:}
\begin{enumerate}
  \item Soil N exists in four pools: NO$_3^-$, NH$_4^+$, organic N, and urea; only NO$_3^-$ and NH$_4^+$ are immediately plant-available
  \item Daily N processes: mineralisation $\to$ nitrification $\to$ plant uptake $\to$ leaching
  \item N response curves follow diminishing returns: large yield gains at low N, small gains at high N
  \item AEN = (Yield$_{N}$ $-$ Yield$_0$) / N rate; typical range 10--16 kg grain/kg N for dryland wheat
  \item N stress reduces LAI, biomass, grain number, and Harvest Index; raises grain protein
  \item Optimal N rate depends on grain price, fertiliser cost, and seasonal rainfall --- use cost-benefit analysis
\end{enumerate}
\end{keyconceptbox}

\section{What's Next?}

Module 5 covers phenology and thermal time (how development depends on temperature).

% ============================================================================
% END MODULE 4
% ============================================================================
